

    \item Two coherent monochromatic point sources \(S_1\) and \(S_2\) of wavelength \(\lambda = 600 \text{ nm}\) are placed symmetrically on either side of the center of the circle as shown. The sources are separated by a distance \(d = 1.8 \text{ mm}\). This arrangement produces interference fringes visible as alternate bright and dark spots on the circumference of the circle. The angular separation between two consecutive bright spots is \(\Delta\Theta\). Which of the following options is/are correct?
        \begin{tasks}(1)
            \task A dark spot will be formed at the point \(P_2\)
            \task At \(P_2\), the order of the fringe will be maximum
            \task The total number of fringes produced between \(P_1\) and \(P_2\) in the first quadrant is close to 3000
            \task The angular separation between two consecutive bright spots decreases as we move from \(P_1\) to \(P_2\) along the first quadrant
        \end{tasks}


\begin{center}
\begin{tikzpicture}
\node {Diagram};
\end{tikzpicture}
\end{center}
